\chapter{Einleitung}
\label{chap:introduction}

\section{Motivation}
Große Sprachmodelle (Large Language Models, LLMs) werden zunehmend in der Softwareentwicklung eingesetzt und unterstützen Entwurf, Implementierung und Wartung.
Im Folgenden wird ein solches \emph{interaktives} Assistenzsystem als \emph{KI-Assistenzsystem} bezeichnet (engl.: \emph{AI Coding Agent}). Der Begriff wird hier im Sinne eines Assistenzsystems ohne eigene Aufgabenplanung verwendet.
In vielen Arbeitsabläufen verschiebt sich der Schwerpunkt von der reinen Codeerzeugung hin zur Auswahl und Integration vorgeschlagener Änderungen.
Damit gewinnt eine Arbeitsweise an Relevanz, die häufig als \emph{Vibe-Coding} beschrieben wird.
Entwicklerinnen und Entwickler delegieren Teilaufgaben an ein KI-Assistenzsystem und übernehmen verstärkt Aufgaben der semantischen Kontrolle, des Reviews sowie der Einbettung in bestehende Strukturen.
Diese Form der Mixed-Initiative-Zusammenarbeit betont weniger syntaktische Details als vielmehr das Verständnis von Anforderungen, Nebenbedingungen und Systemkontext.

Im Brownfield-Kontext ist diese Entwicklung besonders bedeutsam.
Ein Großteil praktischer Softwareentwicklung findet in bestehenden Codebasen statt, die technische Schulden, heterogene Architekturen, historische Entscheidungen und operative Randbedingungen aufweisen.
Bei der Integration von KI-Assistenzsystemen in etablierte Entwicklungsprozesse werden neben Effizienzgewinnen auch Qualitätssicherung, Nachvollziehbarkeit und Risikobegrenzung als relevante Aspekte diskutiert.
Daraus ergibt sich ein Spannungsfeld zwischen schneller Iteration und belastbarem Abgleich.

Für die vorliegende Arbeit ist dabei zentral: Potenzielle Effizienzgewinne werden nicht quantifiziert; im Fokus steht vielmehr, wie Integration, Abgleich und Review im Zusammenspiel mit einem KI-Assistenzsystem empirisch beobachtbar werden.

Wissenschaftlich leistet die Arbeit damit einen Beitrag zur empirischen Auswertung von KI-Assistenzsystemen im Brownfield-Kontext, der über reine Produktivitätsbetrachtungen hinausgeht.

Vor diesem Hintergrund ist der Nutzen eines KI-Assistenzsystems in dieser Arbeit nicht als reine Eigenschaft des Tools zu verstehen, sondern als Zusammenspiel aus (i) Interaktion zwischen Mensch und System, (ii) entwicklerseitiger Kompetenz und (iii) Praktiken des Abgleichs und der Integration.
Automation Bias kann dazu beitragen, dass Vorschläge überbewertet oder unzureichend kontrolliert werden, während umgekehrt übermäßige Skepsis potenzielle Nutzenaspekte reduziert.
Eine wissenschaftlich belastbare Einordnung macht daher eine Betrachtung erforderlich.
Diese verknüpft Kompetenz als Analyseperspektive mit beobachtbaren Artefakten und Ergebnissignalen der Arbeit am Code, ohne daraus kausale Effekte oder Tool-Rangordnungen abzuleiten.

\section{Problemstellung}
In der Diskussion um KI-Assistenzsysteme werden häufig Produktivitätsannahmen formuliert, ohne hinreichend zu differenzieren, wie stark die Interaktion durch Kompetenz, Arbeitsstil und Abgleichspraktiken moderiert wird.
Unter realitätsnahen Bedingungen bestehender Codebasen und unter Zeitdruck bleibt zudem unklar, in welchem Verhältnis wahrgenommener Fortschritt, implementierte Änderungen und als Execution observed ausgewiesenes Verhalten zueinander stehen.
Damit entsteht eine Forschungslücke: Kontrollierte Untersuchungen kombinieren selten Brownfield-Settings, eine differenzierte Ergebnislogik entlang einer Pipeline und Kompetenz als explizite Analyseperspektive.

Diese Arbeit ist daher keine Produktivitätsmessung (z.B. Output-Geschwindigkeit, Codeumfang oder Zeitersparnis), sondern eine Analyse von Integrations-, Abgleichs- und Reviewarbeit im Vibe-Coding-Kontext.

Sie trifft keine Aussagen über kausale Effekte und leitet keine Rangordnung von Tools oder Modellen ab.

Eine zentrale methodische Herausforderung liegt in der Messbarkeit von Zielerreichung.
Chat-Interaktionen und Planäußerungen im Assistenzdialog sind nicht gleichzusetzen mit funktionsfähiger Implementierung.
Ebenso ist eine Implementierung im Code nicht gleichzusetzen mit einem Execution observed-Ausweis dafür, dass gewünschtes Verhalten tatsächlich erreicht wurde.
Diese Differenzen sind im Kontext von KI-generiertem Code besonders relevant, weil die Geschwindigkeit der Generierung und die Plausibilität der Vorschläge die Notwendigkeit eines getrennten Statusausweises nicht aufheben, sondern eher erhöhen können.

Zur Bearbeitung dieser Messprobleme wird eine Diagnoselogik benötigt, die Ergebnisse schrittweise differenziert und durch mehrere Evidenzquellen absichert.
In dieser Arbeit wird dafür eine Pipeline-Logik verwendet, die den Bearbeitungsstand einer Aufgabe als \textit{Attempted}, \textit{Implemented\textsubscript{static}} oder \textit{Execution observed} beschreibt.
Für die Trichterdarstellung wird in Kapitel~6 zusätzlich eine abgeleitete Evidenzstufe  \textit{Implemented}\,$^\ast$ verwendet.
Diese Status sind als Diagnose- und Ergebnislogik zu verstehen, nicht als Qualitätsurteil.
Die Einordnung wird nicht aus einer einzelnen Quelle abgeleitet, sondern trianguliert über Code-Artefakte, Ausführungsartefakte, Chat-Logs und Survey-Daten.
Ziel ist eine methodisch robuste, deskriptive Einordnung, die subjektive oder dialogbasierte Signale nicht mit Execution observed-Ausweisen gleichsetzt und den Interaktionskontext dennoch berücksichtigt.

\section{Zielsetzung der Arbeit}
Ziel dieser Arbeit ist die Auswertung eines KI-Assistenzsystems im Vibe-Coding-Kontext in einer kontrollierten Brownfield-Sandbox.
Zur begrifflichen und methodischen Trennung wird zwischen der Assistenzumgebung GitHub Copilot (Tool), dem zugrundeliegenden LLM Claude Sonnet~4.5 und dem KI-Assistenzsystem als funktionaler Einheit unterschieden.
Beobachtungen beziehen sich je nach Datenquelle auf unterschiedliche Ebenen.
Ein Tool- oder Modellvergleich ist nicht Gegenstand der Arbeit.
Die Wahl von GitHub Copilot erfolgt aus Gründen der ökologischen Validität: Der Integrationskontext (BuyIn) sowie die Zielsetzung der Arbeit orientieren sich an einer Microsoft-nahen Entwicklungsumgebung (insbesondere VS Code und GitHub-Tooling), in der Copilot als praxisnahe Assistenzumgebung naheliegt.
Die Datenerhebung erfolgt in einer kontrollierten Coding-Challenge mit $N = 18$ Teilnehmenden, einer Timebox von 60 Minuten und einer standardisierten, containerisierten Arbeitsumgebung.
Die Aufgaben umfassen acht Tasks (T1--T8).
In der Ergebnisanalyse werden insbesondere frühe Tasks (T1--T5) berücksichtigt, da das Zeitlimit die Bearbeitungsreihenfolge und -tiefe strukturell beeinflusst.

Gegenstand der Arbeit ist dabei ein \emph{interaktives} KI-Assistenzsystem mit kontinuierlicher menschlicher Kontrolle (Human-in-the-Loop); nicht betrachtet werden autonome Agenten mit selbstständiger Aufgabenplanung, Tool-Auswahl oder Entscheidungsautonomie.

Die Arbeit leistet im Kern vier Beiträge:
\begin{enumerate}
    \item[(a)] Entwicklung und Verwendung einer Brownfield-Sandbox als Messumgebung für kontrollierte, aber realitätsnahe Aufgaben in einer bestehenden Codebasis.
    \item[(b)] Operationalisierung von Aufgabenerfolg über die Pipeline-Status: \\
          	extit{Attempted}, \textit{Implemented\textsubscript{static}} und \textit{Execution observed} (sowie \textit{Implemented}\,$^\ast$ als abgeleitete Evidenzstufe) als Trichter- und Diagnoselogik.
    \item[(c)] Analyse der Ergebnisse im Kontext von Kompetenzgruppen als Analyseperspektiven (No-Code, Low-Code, Mid-Code, High-Coder), ohne diese Einteilung als Wertung von Personen zu interpretieren.
    \item[(d)] Methodische Triangulation über Code-Artefakte, Ausführungsartefakte, Chat-Logs und Survey, um die Differenz zwischen intendiertem Vorgehen, implementierten Änderungen und als Execution observed ausgewiesenem Verhalten empirisch einordnen zu können.
\end{enumerate}

\paragraph{Forschungsfragen.}
Aus diesen Zielsetzungen ergeben sich drei Forschungsfragen:
\begin{itemize}
    \item  \textbf{RQ1:} Wie hängt die Entwicklerkompetenz als Analyseperspektive mit dem Erreichen der Pipeline-Status  \textit{Attempted},  \textit{Implemented\textsubscript{static}} und  \textit{Execution observed} (sowie der abgeleiteten Evidenzstufe  \textit{Implemented$^\ast$}) in einer kontrollierten Brownfield-Sandbox zusammen?
    \item  \textbf{RQ2:} Wie tragen Code- und Ausführungsartefakte zur Ergebniseinordnung entlang der Pipeline-Status bei, insbesondere zur Trennung von Implementierung und Execution observed-Ausweisen?
    \item  \textbf{RQ3:} Welche Chat-Interaktionsmuster treten im Vibe-Coding-Kontext auf, und wie sind sie für die Auswertung eines KI-Assistenzsystems im Hinblick auf Abgleich und Integration einzuordnen?
\end{itemize}
Die Kompetenzgruppen dienen dabei ausschließlich der analytischen Typisierung; daraus wird keine Rangordnung abgeleitet und die Einleitung nimmt keine Ergebnisse vorweg.
Die Aussagen dieser Arbeit beziehen sich auf Branches als Analyseeinheiten und auf aggregierte Muster über Gruppen hinweg, nicht auf eine Leistungsbewertung einzelner Personen.
Ergänzend werden theoriegeleitete Perspektiven formuliert, um erwartbare Muster im Spannungsfeld zwischen Unterstützungspotenzial, Abgleichsaufwand und Interaktionsdynamik strukturiert zu beschreiben, ohne Ergebnisse vorwegzunehmen.

\section{Aufbau der Arbeit}
Kapitel~2 führt die theoretischen Grundlagen ein und ordnet Vibe-Coding als Arbeitsstil in relevante Perspektiven ein, darunter Kompetenzmodelle, Human-Computer Interaction, Cognitive Load, Automation Bias sowie etablierte Bezugsrahmen aus Qualitäts- und Produktivitätsdiskursen (u.\,a. ISO/SPACE).
Kapitel~3 beschreibt das Forschungsdesign, die Herleitung theoriegeleiteter Perspektiven und die Operationalisierung zentraler Konstrukte, einschließlich der Kompetenzgruppen als Analyseperspektiven und der Pipeline-Status  \textit{Attempted},  \textit{Implemented\textsubscript{static}} und  \textit{Execution observed} (sowie  \textit{Implemented$^\ast$}).
Kapitel~4 erläutert die Durchführung der Studie, die Datenerhebung und die verwendeten Metriken sowie die Diskussion der Gütekriterien durch die Kombination mehrerer Datenquellen.

Kapitel~5 dokumentiert die technische Brownfield-Sandbox, ihre Architektur und die Automatisierungs- und Auswertungsinfrastruktur, die eine standardisierte, reproduzierbare Messung im kontrollierten Setting ermöglicht.
Kapitel~6 präsentiert die Ergebnisse in deskriptiver Form und strukturiert sie entlang der Pipeline-Logik und der triangulierten Evidenz aus Artefakten und Interaktionen.
Kapitel~7 diskutiert die Befunde, ordnet sie in den theoretischen Kontext ein und fasst Limitationen zusammen, ohne kausale Schlussfolgerungen über die beobachteten Zusammenhänge hinaus zu beanspruchen.
Kapitel~8 schließt die Arbeit mit einer synthetischen Beantwortung der Forschungsfragen und einer Einordnung der Beiträge ab.

\paragraph{Abgrenzung und wissenschaftlicher Beitrag.}
Im Unterschied zu Teilen der bestehenden Forschung steht in dieser Arbeit weder die Optimierung von Prompts noch der Vergleich von Tools oder Modellen im Vordergrund.
Ebenso werden keine Produktivitätsmetriken im Sinne von Ausstoß, Geschwindigkeit oder Zeitersparnis als primäre Zielgrößen herangezogen.
Stattdessen wird Vibe-Coding im Brownfield-Setting als Interaktions- und Integrationsprozess untersucht, der sich nicht zuverlässig über einzelne Proxy-Metriken abbilden lässt.
Der Beitrag liegt in einer auswertenden Perspektive auf Nutzungstypen (Kompetenzgruppen als Analyseperspektiven), einer evidenzbasierten Ergebnislogik über Pipeline-Status sowie der systematischen Triangulation über Code-, Ausführungs-, Chat- und Survey-Artefakte.
Damit wird nachvollziehbar, wie sich intendiertes Vorgehen, implementierte Änderungen und als Execution observed ausgewiesenes Verhalten unter Zeitdruck zueinander verhalten, ohne daraus eine Leistungsbewertung einzelner Personen oder eine Rangordnung von Assistenzsystemen abzuleiten.
